\documentclass[a4paper,12pt]{article}

\usepackage[a4paper,margin=1in]{geometry}
\usepackage{enumitem}
\usepackage{listings}
\usepackage{xcolor}
\usepackage{hyperref}

\setlength{\parindent}{0pt}
\setlength{\parskip}{0.6em}

% Assembly style
\lstdefinestyle{AsmStyle}{
    language={[x86masm]Assembler},
    basicstyle=\ttfamily\footnotesize,
    keywordstyle=\color{blue},
    commentstyle=\color{green!40!black},
    numbers=left,
    numberstyle=\tiny,
    breaklines=true,
    frame=single
}

\title{Project 2 Report: \texttt{schedsim}\\
\large CMPE 230 Systems Programming -- Spring 2026}

\author{
Student 1 Name, Student 1 ID \\
Student 2 Name, Student 2 ID
}

\date{\today}

\begin{document}

\maketitle

\section*{Report Guidelines}

This report must explain \textbf{how your assembly program works internally}.

Avoid generic statements. Focus on:
\begin{itemize}
\item how parsing is implemented at register level,
\item how program state is stored in memory,
\item how scheduling decisions are computed,
\item how control flow is structured in assembly.
\end{itemize}


\section{Project Overview}

Briefly describe:
\begin{itemize}
\item what your scheduler does,
\item supported algorithms (FCFS, SJF, SRTF, PF, RR),
\item overall execution flow from input to output.
\end{itemize}


\section{Parsing Logic}

\subsection{Input Processing Strategy}

Explain how you parse the single input line:

\begin{itemize}
\item how tokens are separated (spaces, hyphens),
\item how algorithm type is detected,
\item how numeric values are parsed from characters.
\end{itemize}

\subsection{Register-Level Parsing Design}

Describe explicitly:

\begin{itemize}
\item which registers hold pointers (e.g., input buffer),
\item which registers hold temporary numeric values,
\item how digits are accumulated into integers,
\item how delimiters (' ', '-') are handled.
\end{itemize}

Explain how your parsing avoids corruption and maintains correctness.

\section{State Representation (Memory Layout)}

Explain how all scheduling data is stored.

\subsection{Process Representation}

Describe how each process is represented:

\begin{itemize}
\item ID (character),
\item burst time,
\item arrival time,
\item remaining time,
\item priority (if applicable).
\end{itemize}

Explain whether you used:
\begin{itemize}
\item arrays,
\item structs (manual layout),
\item separate parallel arrays.
\end{itemize}

\subsection{Global State}

Explain:
\begin{itemize}
\item current time,
\item ready queue / process list,
\item output buffer,
\item counters (remaining processes).
\end{itemize}


\section{Register Usage Strategy}

Clearly document your register design:

\begin{itemize}
\item which registers are used for:
    \begin{itemize}
    \item loop counters,
    \item pointers,
    \item current process index,
    \item temporary computations,
    \end{itemize}
\item which registers are preserved across function calls (if any),
\item how you avoid accidental overwrites.
\end{itemize}

If you used conventions (caller-saved / callee-saved), explain them.

\section{Scheduling Logic}

\subsection{Clock Cycle Simulation}

Explain how one time step is executed:

\begin{itemize}
\item how eligible processes are determined,
\item how the next process is selected,
\item how output is appended,
\item how state is updated.
\end{itemize}

\subsection{Algorithm Implementations}

Briefly explain each:

\begin{itemize}
\item FCFS: arrival-based ordering
\item SJF: sorting or selection by burst
\item SRTF: per-cycle minimum remaining time
\item PF: priority + tie-breaking
\item RR: queue rotation + quantum handling
\end{itemize}

Focus on differences in control flow and data updates.

\section{Control Flow Design}

Describe how your assembly program is structured:

\begin{itemize}
\item main loop (clock cycle loop),
\item parsing phase,
\item scheduling phase,
\item output phase,
\item termination condition.
\end{itemize}

Mention use of:
\begin{itemize}
\item labels,
\item jumps (conditional / unconditional),
\item loop patterns.
\end{itemize}


\section{Representative Code Snippet}

\begin{lstlisting}[style=AsmStyle]
# Example snippet
\end{lstlisting}

Explain:
\begin{itemize}
\item what the code does,
\item which registers are used,
\item how control flow works,
\item why it is correct.
\end{itemize}

\section{Testing and Debugging}

\subsection{Testing Strategy}

\begin{itemize}
\item sample inputs,
\item edge cases (ties, idle CPU, preemption),
\item manual vs automated testing.
\end{itemize}

\subsection{Debugging Approach}

\begin{itemize}
\item print-based debugging,
\item register inspection,
\item step-by-step tracing.
\end{itemize}

\subsection{Edge Cases}

Examples:

\begin{itemize}
\item multiple arrivals at same time,
\item identical priorities,
\item quantum partially used (RR),
\item CPU idle periods.
\end{itemize}

\section{Course Concepts}

Relate your implementation to course topics:

\subsection{Memory and Data Layout}

How memory is organized and accessed.

\subsection{Registers and Low-Level Computation}

How computation is performed without high-level abstractions.

\subsection{Control Flow and Branching}

How loops and decisions are implemented.

\subsection{Parsing Without Libraries}

Manual string processing in assembly.

\subsection{Tradeoffs}

\begin{itemize}
\item simplicity vs efficiency,
\item readability vs performance,
\item static vs dynamic design.
\end{itemize}

\section{Conclusion}

Summarize:

\begin{itemize}
\item strongest design decision,
\item biggest challenge,
\item possible improvements.
\end{itemize}

\section*{AI Usage Statement}

The report must be written by you.

If AI tools were used:
\begin{itemize}
\item which tools,
\item for what purpose,
\item confirmation that all technical content is understood.
\end{itemize}

\end{document}